\subsection{Historický vývoj}
\label{subsec:llm-history}

Dnešní LLM jsou výsledkem několika desetiletí výzkumu v oblasti zpracování přirozeného jazyka, neuronových sítí a strojového učení.
Pochopení této vývojové linie pomáhá vysvětlit, proč jsou dnešní modely takové, jaké jsou, a jaké problémy se přitom řešily.

\subsubsection*{Pravidlové systémy (1960 -- 1980)}
\label{subsubsec:llm-history-rules}

Prvními systémy zpracování přirozeného jazyka byly programy pracující na bázi ručně definovaných pravidel a vzorů.
Nejznámějším příkladem je chatbot ELIZA z roku 1966, který simuloval psychoterapeuta pomocí jednoduchých vzorů přepisování vstupních vět.
ELIZA přes svou primitivnost vyvolávala v uživatelích dojem skutečného porozumění, jev, který Joseph Weizenbaum nazval \textit{ELIZA effect}.
Pravidlové systémy měly však zásadní omezení.
Každá nová situace vyžadovala nová explicitní pravidla a systém nedokázal generalizovat na neočekávané vstupy~\cite{weizenbaum-eliza}.

\subsubsection*{Statistické jazykové modely a n-gramy (1980 -- 2000)}
\label{subsubsec:llm-history-statistical}

Přechod od pravidel ke statistickým přístupům přinesl výrazný posun.
N-gramové modely odhadovaly pravděpodobnost výskytu slova na základě $n-1$ předcházejících slov, přičemž pravděpodobnosti byly odhadována z četností v trénovaným korpusu.
Bigramový model například odhaduje $P(w_t \mid w_{t-1})$ a trigramový model rozšiřuje podmínění o jeden token zpět.

Nevýhoda je, že kontextové okno bylo omezeno na $n$ tokenů, přičemž při větším $n$ exponenciálně rostla potřeba dat pro spolehlivý odhad pravděpodobností.
Model tak nemohl zachytit dlouhodobé závislosti v textu.
Například to, čemu odkazuje zájmeno na konci věty nebo jaké téma se rozvíjí v delším odstavci.

\subsubsection*{Neuronové jazykové modely a word embeddingy (2000 -- 2017)}
\label{subsubsec:llm-history-neural}

Nástup neuronových sítí přinesl dvě klíčové inovace.
První z nich jsou \textit{word embeddingy}, husté vektorové reprezentace slov, ve kterých slova s podobným významem leží geometricky blízko sebe.
Modely Word2Vec (2013)~\cite{mikolov-word2vec} a GloVe (2014)~\cite{pennington-glove} ukázaly, že lze natrénovat reprezentace zachycující sémantické vztahy: král, můž, žena se blíží slovu královna.

Druhou inovací jsou rekurentní neuronové sítě (RNN) a zejména jejich varianta LSTM (Long Short-Term Memory)~\cite{hochreiter-lstm}.
LSTM zavedlo výpočetní brány, které modelu umožňují selektivně uchovávat informace z předchozích kroků, čímž byl z velké části vyřešen problém mizejícího gradientu (vanishing gradient), který brzdil trénování hlubokých rekurentních sítí.
Přesto měly RNN a LSTM dvě přetrvávající omezení: zpracovávaly tokeny sekvenčně, což znemožňovalo efektivní paralelizaci, a v praxi měly obtíže se zachycením závislostí na velmi vzdálených tokenech.

\subsubsection*{Architektura Transformer (2017)}
\label{subsubsec:llm-history-transformer}

Zásadní průlom nastal v roce 2017 s publikací práce \textit{Attention is All You Need} od výzkumníků z Google Brain~\cite{vaswani-attention}.
Transformer zcela nahradil rekurentní zpracování mechanismem \textit{self-attention}, který umožňuje modelu uvažovat o vztazích mezi \textit{všemi} tokeny vstupu najednou v jednom paralelním výpočtu.
Tím byly odstraněny obě hlavní nevýhody RNN: výpočty bylo možné plně paralelizovat na GPU a model mohl přirozeně zachytit závislosti mezi libovolně vzdálenými tokeny.

\subsubsection*{Období předtrénovaných modelů a škálování (2018 -- současnost)}
\label{subsubsec:llm-history-pretraining}

Demonstrace síly předtrénování, kdy je model natrénovaný na velkém korpusu a následně doladěný na konkrétní úlohu dosahoval výsledků překonávajících dosavadní specializované přístupy.
GPT-3 (2020, 175 miliard parametrů) ukázal, že se škálováním modelů objevují kvalitativně nové schopnosti~\cite{brown-gpt3}.
Od roku 2022 nastupuje éra instrukčního dolaďování a modelů trénovaných zpětnou vazbou od lidí (RLHF), jejichž výsledkem jsou aktuální asistenti jako ChatGPT, Claude nebo Geminy.

\begin{table}[H]
    \centering
    \resizebox{\textwidth}{!}{%
        \begin{tabular}{lll}
            \toprule
            \textbf{Období}   & \textbf{Klíčová technologie}    & \textbf{Zástupce}         \\
            \midrule
            1960--1980        & Pravidlové systémy              & ELIZA (1966)~\cite{weizenbaum-eliza}              \\
            1980--2000        & N-gramové modely                & IBM Candide               \\
            2000--2013        & Neuronové sítě, word embeddingy & Word2Vec (2013)~\cite{mikolov-word2vec}           \\
            2013--2017        & RNN, LSTM                       & seq2seq (2014)            \\
            2017--2020        & Transformer, předtrénování      & BERT (2018)~\cite{devlin-bert}, GPT-2 (2019) \\
            2020--2022        & Škálování, few-shot learning    & GPT-3 (2020)~\cite{brown-gpt3}              \\
            2022--dnes        & Instrukční ladění, RLHF         & ChatGPT, Claude, Llama    \\
            \bottomrule
        \end{tabular}
    }
    \caption{Přehled klíčových etap vývoje jazykových modelů}
    \label{tab:llm-history-summary}
\end{table}

\endinput